\chapter{Adaptive Defense Policy Network (\adpn{})}
\label{ch:adpn}

\section{Problem Formulation}

We model the defense problem as an MDP
$(\mathcal{S}, \mathcal{A}, \mathcal{R}, \mathcal{T}, \gamma)$:
\begin{itemize}
  \item $\mathcal{S} \subset \mathbb{R}^{16}$: network state space
  \item $\mathcal{A} = \{0,\ldots,7\}$: 8 discrete defense actions
  \item $\mathcal{R}(s,a) \in [-1.5, +1.0]$: defense reward minus FPR cost
  \item $\mathcal{T}$: transitions induced by the \caesar{} environment
  \item $\gamma = 0.95$: discount factor
\end{itemize}

\section{Dueling Double-DQN Architecture}

\subsection{Dueling Streams}
\begin{equation}
  Q(s,a;\theta,\alpha,\beta) = V(s;\theta,\beta)
    + A(s,a;\theta,\alpha)
    - \frac{1}{|\mathcal{A}|}\sum_{a'} A(s,a';\theta,\alpha)
\end{equation}

Shared feature extractor: Linear(16,128)--ReLU--Linear(128,128)--ReLU.
Value stream: Linear(128,64)--ReLU--Linear(64,1).
Advantage stream: Linear(128,64)--ReLU--Linear(64,8).

\subsection{Double DQN Target}
Decoupling action selection from evaluation:
\begin{equation}
  y_t = r_t + \gamma\, Q_{\hat\theta}\!\left(s_{t+1},
    \argmax_{a'} Q_\theta(s_{t+1},a')\right)
\end{equation}
This reduces the overestimation bias that standard DQN exhibits when
$Q$-values are noisy — particularly important in the early co-evolutionary
phase when attack patterns are highly non-stationary.

\section{Replay Buffer and Sampling}

A circular replay buffer of capacity 12,000 stores tuples
$(s_t, a_t, r_t, s_{t+1}, \mathrm{done})$.
Mini-batches of size 64 are sampled uniformly.
The target network $Q_{\hat\theta}$ is synchronized with $Q_\theta$
every 100 gradient steps.

\section{Proactive Override Mechanism}

When the \tagraph{} produces a proactive defense recommendation $p_t$
with confidence $\geq \tau_{\mathrm{pro}} = 0.55$, \adpn{} applies
$a_t = p_t$ directly, bypassing the $\varepsilon$-greedy policy.
This transition is stored in the replay buffer as a standard
$(s_t, p_t, r_t, s_{t+1})$ tuple so the network learns to associate
the predicted attack context with the optimal countermeasure.

\section{Co-evolutionary Reward Shaping}

The environment reward $r_d^{\mathrm{env}}$ is shaped with a
false-positive penalty:
\begin{equation}
  r_t = r_d^{\mathrm{env}}(a_t) - 0.4 \cdot \mathrm{FPR}(a_t)
\end{equation}
This prevents the agent from over-deploying aggressive actions
(e.g.\ node isolation) when a milder response (rate limiting) suffices.

\section{Convergence and $\varepsilon$-Decay}

$\varepsilon$ decays exponentially: $\varepsilon_{t+1} = \max(0.05,
0.995 \cdot \varepsilon_t)$, reaching minimum in approximately
600 steps. The agent transitions from exploration-dominant to
exploitation-dominant behaviour around episode 60 of 150, which
coincides with the \tagraph{} accumulating sufficient transition
statistics for reliable proactive predictions.
